% ============================================================
%  defs.tex — comandes i entorns propis
% ============================================================

% ---------- Textos generals ----------
\newcommand{\titoldoc}{Errades del solucionari}
\newcommand{\curs}{Matemàtiques II \textperiodcentered{} 2n de Batxillerat}

% ---------- Notació ----------
\newcommand{\R}{\mathbb{R}}
\newcommand{\dlim}{\displaystyle\lim}   % límit amb el subíndex a sota

% ---------- Colors ----------
\colorlet{fonscaixa}{blue!4}
\colorlet{marccaixa}{gray!100!black}
\colorlet{colorondiu}{red!65!black}
\colorlet{colorhadedir}{green!40!black}

% ---------- Capçalera de cada setmana ----------
%   \setmana{número}{unitat}{data límit}{exercicis de la setmana}
\newcommand{\setmana}[4]{%
  \section*{Setmana #1}%
  \addcontentsline{toc}{section}{Setmana #1}%
  \noindent\textbf{#2}\quad\textbullet\quad data límit: #3\par
  \noindent Exercicis de la setmana: #4.\par
  \medskip}

% ---------- Una errada ----------
%   \begin{errada}{Exercici 94 b)}
%     \ondiu{el que hi ha imprès al solucionari}
%     \hadedir{el que hi hauria de dir}
%     \motiu{explicació breu}
%   \end{errada}
% El «Motiu» només surt al PDF si la flag de dalt de main.tex és \motiutrue;
% amb \motiufalse, el text de \motiu{...} s'ignora.
\newenvironment{errada}[1]{%
  \begin{tcolorbox}[colback=fonscaixa,colframe=marccaixa,
    title={#1},fonttitle=\bfseries]
  \begin{description}[style=multiline,leftmargin=2.3cm,
    font=\normalfont\bfseries,itemsep=3pt,topsep=0pt,parsep=0pt]}%
 {\end{description}\end{tcolorbox}}

\newcommand{\ondiu}[1]{\item[{\color{colorondiu}On diu}] #1}
\newcommand{\hadedir}[1]{\item[{\color{colorhadedir}Ha de dir}] #1}
\newcommand{\motiu}[1]{\ifmotiu\item[{\color{marccaixa}Motiu}] #1\fi}

% ---------- Una observació (no és cap errada, però pot confondre) ----------
%   \begin{observacio}{Exercici 120} ... \end{observacio}
% Al fitxer de cada setmana, el bloc d'observacions va entre \ifobs i \fi:
% només surt al PDF si la flag de dalt de main.tex és \obstrue.
\newtcolorbox{observacio}[1]{colback=gray!6,colframe=gray!60!black,
  title={Observació \textperiodcentered{} #1},fonttitle=\bfseries}

% ---------- Exercicis revisats i sense errades ----------
%   \senseerrades{103, 40 i 106}
% Als fitxers de les setmanes aquesta línia hi és escrita però comentada (%):
% per fer-la sortir al PDF, treu-li el %.
\newcommand{\senseerrades}[1]{%
  \medskip\noindent$\checkmark$\ \ \textbf{Revisats i sense errades:} #1.\par}
